\section{Geometric contacts within one cell fan}
\label{sec:area_law_fan_run_contacts}

The merging of triangles in \cite[Section~11]{OpenAI2026AreaLaw} uses their
geometric contacts. The following local assertions hold for every origin,
dyadic exponent, signed cell index and optional side-midpoint subdivision.
They require no assignment of labels outside the cell.

% Source: 10-geometry.tex, prop:two-families, lines 308--323,
% revision adc7f1241b42e322a6451854ab7e4b4c146bf78a.
% These are within-cell assertions, not the full two-family proposition.

\begin{theorem}[Contact determines consecutive triangles]
    \label{thm:area_law_fan_triangle_contact}
    \lean{TNLean.PEPS.AreaLaw.Geometry.cellFanPolygons_nontrivial_inter_cases}
    \leanok
    \uses{def:area_law_cell_fan}
    Fix a fan with center $c$ in a translated dyadic square. Write
    $e_i=[a_i,b_i]$ for its elementary perimeter segments, oriented
    counterclockwise, and $P_i=\operatorname{conv}\{c,a_i,b_i\}$ for its
    triangles. If $i\ne j$ and $P_i\cap P_j$ contains two distinct points,
    then one of the following alternatives holds:
    \begin{align}
        b_i & =a_j, & P_i\cap P_j & =[c,b_i], \\
        b_j & =a_i, & P_i\cap P_j & =[c,b_j].
    \end{align}
    Thus geometric contact determines consecutive perimeter segments and
    identifies the complete common radial edge.
\end{theorem}
\begin{proof}\leanok
    \uses{thm:area_law_cell_fan_intersections,thm:area_law_unsplit_side_endpoints}
    If the two perimeter segments were disjoint, the fan intersection
    identity would give $P_i\cap P_j=\{c\}$. Their intersection is therefore
    nonempty. Two distinct elementary segments on the same side are its
    first and second halves, whose intersection is the midpoint. Elementary
    segments on distinct sides can meet only at a common corner of
    consecutive sides; any segment containing that corner has it as the
    corresponding endpoint. In either case their intersection is the
    singleton of a last-to-first endpoint. Substituting into the fan
    intersection identity gives the asserted radial segment.
\end{proof}

\begin{theorem}[Opposite labels at contacts between distinct runs]
    \label{thm:area_law_fan_run_geometric_contact}
    \lean{TNLean.PEPS.AreaLaw.Geometry.cellFanRunRegions_contact_colors_ne}
    \leanok
    \uses{def:area_law_fan_runs}
    Give the triangles of this fan arbitrary labels in $\{0,1\}$.
    Let $R\ne T$ be two actual run identifiers, with regions $T_R,T_T$.
    If $T_R\cap T_T$ contains two distinct points, then every triangle
    in $R$ has a different label from every triangle in $T$.
\end{theorem}
\begin{proof}\leanok
    \uses{thm:area_law_fan_triangle_contact,thm:area_law_fan_run_opposite,
        thm:area_law_fan_run_color}
    Choose a common point $x\ne c$. It belongs to triangles $P_i,P_j$
    in the two runs. Their slots are distinct, and both triangles also
    contain $c$, so their intersection contains two distinct points.
    The preceding theorem gives endpoint adjacency. Distinct runs then
    force the two labels to differ. Label constancy on each run gives
    the assertion for every pair of constituent triangles.
\end{proof}

These assertions concern a single fan, whose triangles share one center.
Across different cells, a shared nondegenerate segment must be distinguished
from isolated common corners. The classification of global interfaces is a
further part of the two-family construction.
